\chapter{Introducción}

Las pequeñas y medianas empresas de transporte coordinan cada día personas, vehículos, clientes y envíos. Cuando la información queda repartida entre hojas de cálculo, mensajes y llamadas, una pregunta sencilla ---qué vehículo está libre o qué ocurrió con una entrega--- requiere reconstruir datos de varias fuentes. La fragmentación también favorece errores con impacto directo: asignar el mismo recurso dos veces, emplear datos ya obsoletos o perder la autoría de una incidencia.

TransitOps plantea una aplicación web única para ese trabajo operativo. Permite mantener los catálogos, planificar envíos, asignar recursos, controlar el ciclo de vida y consultar una cronología verificable. Una administración pequeña de usuarios y un resumen de actividad completan el núcleo necesario para que el producto sea demostrable por sí mismo y no una colección de endpoints aislados.

La solución separa una SPA React/TypeScript, una API ASP.NET Core y PostgreSQL, y ofrece ejecución reproducible mediante contenedores. El desarrollo se organiza en rebanadas verticales: primero un esqueleto autenticado y después catálogos, envíos, operación, trazabilidad y administración. Los seis primeros sprints cubren RF-01 a RF-14, y el séptimo añade el endurecimiento, las pruebas de sistema y el despliegue accesible.

\section{Contexto del trabajo}

El trabajo se desarrolla como TFG del Grado en Ingeniería Informática de la Universidade da Coruña, en la mención de Ingeniería del Software. Su objeto es el diseño y la construcción íntegra de una aplicación de gestión operativa para una empresa de transporte, junto con el ciclo de vida de ingeniería que la produce: desde la elicitación de necesidades hasta el despliegue verificado del producto.

La situación de partida combina dos necesidades. Desde el dominio, una empresa de transporte pequeña requiere visibilidad compartida y reglas que eviten inconsistencias. Desde el TFG, se necesita demostrar competencias de ingeniería más amplias que la programación de una capa: comprender a la parte interesada, delimitar alcance, justificar arquitectura, probar comportamiento, reproducir entornos y evaluar el producto. TransitOps une ambos objetivos en un caso suficientemente rico para contener estados, relaciones, roles y trazabilidad, pero lo bastante pequeño para completarlo.

El interés del proyecto no reside, por tanto, únicamente en implementar funciones. El objetivo académico es aplicar y documentar de forma disciplinada el ciclo de vida completo del software: elicitación, especificación, diseño, implementación, verificación, despliegue y análisis. Cada decisión se relaciona con una necesidad, se materializa en backend y frontend cuando corresponde y se acompaña de pruebas y evidencia. Esta trazabilidad permite valorar no solo qué hace la aplicación, sino cómo se llegó de un problema expresado en lenguaje de negocio a un sistema mantenible.

\section{Objetivos}

El objetivo general es diseñar, construir, probar y desplegar una aplicación web mantenible que permita gestionar las operaciones esenciales de una empresa de transporte y que constituya una demostración verificable del ciclo de vida completo del software. De ese objetivo se derivan los siguientes objetivos específicos:

\begin{itemize}
  \item Obtener necesidades mediante una entrevista simulada con una responsable del dominio y formalizarlas sin introducir prematuramente decisiones técnicas.
  \item Mantener trazabilidad entre entrevista, RF-01 a RF-14, RN-01 a RN-17, flujos de negocio, diseño, implementación y pruebas.
  \item Diseñar un dominio coherente, un contrato HTTP uniforme y una arquitectura separada de backend, frontend y persistencia.
  \item Implementar incrementos verticales utilizables, con autenticación y autorización desde el comienzo, hasta cubrir el alcance funcional completo.
  \item Automatizar pruebas de reglas en backend, contratos HTTP y comportamiento observable en frontend.
  \item Garantizar ejecución local reproducible mediante Docker Compose, configuración externa y migraciones versionadas, sin secretos reales en el repositorio.
  \item Desplegar el producto en un entorno accesible, ejecutar los cuatro flujos de sistema y analizar los resultados y limitaciones.
  \item Elaborar documentación técnica y académica que permita reproducir decisiones y defender el proceso seguido.
\end{itemize}

\section{Alcance y criterios de éxito}

El alcance se mantiene deliberadamente acotado. Incluye usuarios internos con roles de operador y administrador; vehículos, conductores y clientes; envíos con filtros; asignación conjunta; estados; eventos e incidencias; y un resumen operativo. Quedan fuera la optimización de rutas, el seguimiento GPS, la facturación, la aplicación móvil y el acceso de clientes externos. La exclusión no implica que carezcan de valor: evita que capacidades periféricas impidan cerrar y evaluar de extremo a extremo el núcleo elegido. Tampoco se persigue profundidad en infraestructura, porque la integración continua, los contenedores y el despliegue son medios para reproducir y evaluar el producto; la aportación principal es la disciplina aplicada de ingeniería del software a todo su recorrido.

Los criterios funcionales son el cumplimiento verificable de RF-01 a RF-14 y la ejecución satisfactoria de los cuatro flujos de negocio. En calidad, se exige una compilación reproducible, suites automatizadas correctas, esquema migrable desde cero, autorización efectiva en servidor, errores comprensibles y conservación del historial. En operación, la aplicación debe poder levantarse siguiendo documentación explícita y sin incorporar credenciales reales a ficheros versionados.

Los criterios funcionales se alcanzaron durante los seis primeros incrementos con pruebas de componente e integración y verificaciones manuales por rebanada. El séptimo añadió las pruebas \gls{e2e} de los cuatro flujos, las garantías de concurrencia sostenidas por el motor relacional, una sesión endurecida y el despliegue en un entorno accesible con su evidencia recogida. El éxito solo se declara en ese punto, porque una cobertura amplia y automatizada no equivale a un cierre de sistema demostrado sobre el entorno real.

\section{Estructura de la memoria}

La memoria sigue las fases del ciclo de vida y añade una crónica iterativa. Tras este capítulo expone el marco tecnológico, la metodología ---con la planificación temporal y la estimación de costes--- y los requisitos, y presenta a continuación la arquitectura y el diseño detallado. Los capítulos posteriores conservan la evolución por sprint, la validación y el despliegue; a continuación, la matriz de trazabilidad y los procedimientos de reproducción cierran la cadena de evidencia, y los dos últimos recogen los resultados y las conclusiones. Esta doble lectura separa el diseño resultante de la evidencia temporal que explica cómo se construyó.
